\section{Sucesiones}

Una \emph{sucesión} de números reales es una función \(x\colon \N \to \R\). Se escribe \(x(n) = x_n\) y se lo llama el \(n\)-ésimo término de la sucesión. La sucesión se indica \({(x_n)}_{n \in \N}\) o \((x_1, x_2, \dots)\).

Una \emph{subsucesión} de \(x\) es la restricción de \(x\) a un subconjunto infinito \(A = \{ n_1 < n_2 < \cdots \} \subset \N\). Se escribe \({(x_n)}_{n \in A}\) o \({(x_{n_i})}_{i \in \N}\).

\begin{note}
  Estrictamente, la subsucesión no está definida en \(\N\), pero se la puede considerar como una función definida en \(\N\) componiendo con la biyección \(i \mapsto n_i\). Por eso se usa la notación \({(x_{n_i})}_{i \in \N}\).
\end{note}

\begin{definition}[Límite]
  Decimos que \(a = \lim_{n \to \infty} x_n\) si
  \[
    (\forall \e > 0)(\exists n_0 \in \N) \colon |x_n - a| < \e \quad \forall n > n_0.
  \]
  Equivalentemente, si para todo \(\e > 0\) el intervalo \((a - \e, a + \e)\) contiene a todos los términos de la sucesión salvo, quizás, una cantidad finita.
\end{definition}

\begin{theorem}[Unicidad del límite]
  Si \(\lim_{n \to \infty} x_n = a\) y \(\lim_{n \to \infty} x_n = b\), entonces \(a = b\).
\end{theorem}

\begin{proof}
  Sea \(\e > 0\). Existen \(n_0, n_1 \in \N\) tales que \(|x_n - a| < \e\) para \(n > n_0\) y \(|x_n - b| < \e\) para \(n > n_1\). Luego, para \(n > \max\{n_0, n_1\}\),
  \[
    |a - b| \leq |a - x_n| + |x_n - b| < 2\e.
  \]
  Como \(\e > 0\) es arbitrario, \(|a - b| = 0\), es decir, \(a = b\).
\end{proof}

\begin{theorem}
  Si \(\lim_{n \to \infty} x_n = a\), toda subsucesión de \({(x_n)}_{n \in \N}\) converge a \(a\).
\end{theorem}

\begin{proof}
  Sea \({(x_{n_i})}_{i \in \N}\) una subsucesión. Dado \(\e > 0\), existe \(n_0\) tal que \(|x_n - a| < \e\) para \(n > n_0\). Como \(n_1 < n_2 < \cdots\) es estrictamente creciente, existe \(i_0\) con \(n_{i_0} > n_0\), y entonces para \(i > i_0\) se tiene \(n_i > n_0\) y por lo tanto \(|x_{n_i} - a| < \e\).
\end{proof}

\begin{theorem}
  Toda sucesión convergente es acotada.
\end{theorem}

\begin{proof}
  Sea \(a = \lim_{n \to \infty} x_n\). Tomando \(\e = 1\), existe \(n_1\) tal que \(x_n \in (a - 1, a + 1)\) para todo \(n > n_1\). Sean \(A = \{ x_1, \dots, x_{n_1}, a - 1, a + 1 \}\), \(c = \min A\) y \(d = \max A\). Entonces \(x_n \in [c, d]\) para todo \(n \in \N\).
\end{proof}

\begin{theorem}
  Toda sucesión monótona y acotada es convergente.
\end{theorem}

\begin{proof}
  Supongamos que \({(x_n)}\) es creciente y acotada, y sea \(a = \sup \{ x_n : n \in \N \}\). Dado \(\e > 0\), \(a - \e\) no es cota superior de \(\{ x_n \}\), luego existe \(n_0\) con \(x_{n_0} > a - \e\). Como la sucesión es creciente, para \(n > n_0\) vale
  \[
    a - \e < x_{n_0} \leq x_n \leq a < a + \e,
  \]
  es decir, \(x_n \to a\). El caso decreciente es análogo, con el ínfimo.
\end{proof}

\begin{corollary}
  Si una sucesión monótona tiene una subsucesión convergente, entonces es convergente.
\end{corollary}

\begin{proof}
  La subsucesión es acotada, y como la sucesión es monótona, está acotada por los mismos valores (el conjunto de sus términos está entre \(x_1\) y una cota de la subsucesión). Por el teorema anterior, es convergente.
\end{proof}

\begin{eg}
  Sea \(x_n = a^n\) con \(a \in \R\).
  \begin{itemize}
    \item Si \(a = 0\) o \(a = 1\), la sucesión es constante, luego convergente.
    \item Si \(a = -1\), la sucesión diverge, pues \(x_{2n} = 1 \to 1\) y \(x_{2n+1} = -1 \to -1\) son subsucesiones con límites distintos.
    \item Si \(|a| > 1\), escribimos \(|a| = 1 + h\) con \(h > 0\); por la desigualdad de Bernoulli, \(|a|^n \geq 1 + nh\), y la sucesión no es acotada, luego diverge.
    \item Si \(0 < a < 1\), la sucesión es decreciente y acotada inferiormente por \(0\), luego converge a algún \(l\). Como \(x_{n+1} = a \, x_n\), pasando al límite \(l = a l\), de donde \(l = 0\).
    \item Si \(-1 < a < 0\), se tiene \(|a^n - 0| = |a|^n \to 0\) por el caso anterior, de modo que \(a^n \to 0\).
  \end{itemize}
\end{eg}

\section{Propiedades de los límites}

\begin{theorem}
  Si \(\lim_{n \to \infty} x_n = 0\) e \({(y_n)}_{n \in \N}\) es una sucesión acotada, entonces \(\lim_{n \to \infty} x_n y_n = 0\).
\end{theorem}

\begin{proof}
  Sea \(c > 0\) con \(|y_n| < c\) para todo \(n\). Dado \(\e > 0\), existe \(n_0\) tal que \(|x_n| < \frac{\e}{c}\) para \(n > n_0\). Entonces \(|x_n y_n| < c \cdot \frac{\e}{c} = \e\) para todo \(n > n_0\).
\end{proof}

Sean \(\lim_{n \to \infty} x_n = a\) y \(\lim_{n \to \infty} y_n = b\).

\begin{prop}
  \(x_n + y_n \to a + b\).
\end{prop}

\begin{proof}
  Dado \(\e > 0\), existen \(n_0\) y \(n_1\) tales que \(|x_n - a| < \frac{\e}{2}\) para \(n > n_0\) y \(|y_n - b| < \frac{\e}{2}\) para \(n > n_1\). Para \(n > \max\{n_0, n_1\}\),
  \[
    |(x_n + y_n) - (a + b)| \leq |x_n - a| + |y_n - b| < \e. \qedhere
  \]
\end{proof}

\begin{prop}
  \(x_n - y_n \to a - b\).
\end{prop}

\begin{proof}
  Análoga a la anterior.
\end{proof}

\begin{prop}
  \(x_n y_n \to a b\).
\end{prop}

\begin{proof}
  Se tiene
  \[
    x_n y_n - a b = x_n (y_n - b) + b (x_n - a).
  \]
  Como \(y_n - b \to 0\) y \((x_n)\) es acotada por ser convergente, \(x_n (y_n - b) \to 0\). Además \(b (x_n - a) \to 0\). Luego \(x_n y_n - ab \to 0\).
\end{proof}

\begin{prop}
  Si \(b \neq 0\), entonces \(y_n \neq 0\) salvo finitos términos y \(\dfrac{x_n}{y_n} \to \dfrac{a}{b}\).
\end{prop}

\begin{proof}
  Tomando \(\e = \frac{|b|}{2}\), existe \(n_0\) tal que \(|y_n - b| < \frac{|b|}{2}\) para \(n > n_0\), de modo que \(|y_n| > \frac{|b|}{2} > 0\) para \(n > n_0\); en particular \(y_n \neq 0\). Para \(n > n_0\),
  \[
    \frac{x_n}{y_n} - \frac{a}{b} = \frac{b x_n - a y_n}{b\, y_n} = \frac{1}{b\, y_n} \bigl( b (x_n - a) - a (y_n - b) \bigr).
  \]
  La sucesión \(\bigl( \frac{1}{b y_n} \bigr)_{n > n_0}\) es acotada, pues \(\left| \frac{1}{b y_n} \right| < \frac{2}{b^2}\), y el factor entre paréntesis tiende a \(0\). Por el teorema anterior, \(\frac{x_n}{y_n} - \frac{a}{b} \to 0\).
\end{proof}

\section{Ejemplos con subsucesiones}

\begin{eg}
  Si \(a > 0\), entonces \(\sqrt[n]{a} \to 1\).
\end{eg}

\begin{proof}
  Si \(a = 1\) es trivial. Si \(a > 1\), la sucesión \(x_n = a^{1/n}\) es decreciente: \(a^{1/n} > a^{1/(n+1)} \iff a^{n+1} > a^n\), lo cual es cierto pues \(a > 1\) (se elevó a la \(n(n+1)\)). Además \(x_n > 1\), de modo que converge a un límite \(l \geq 1\). Si \(0 < a < 1\), la sucesión es creciente y acotada por \(1\), y converge a un \(l \geq a > 0\).

  En ambos casos \(l > 0\). La subsucesión \(y_n = a^{\frac{1}{n(n+1)}}\) (índices \(n(n+1)\)) converge también a \(l\), y
  \[
    y_n = a^{\frac{1}{n} - \frac{1}{n+1}} = \frac{a^{1/n}}{a^{1/(n+1)}} \longrightarrow \frac{l}{l} = 1,
  \]
  usando que \(l \neq 0\). Por unicidad del límite, \(l = 1\).
\end{proof}

\begin{eg}
  \(\sqrt[n]{n} \to 1\).
\end{eg}

\begin{proof}
  Veamos que la sucesión es decreciente a partir del tercer término. Elevando a la \(n(n+1)\),
  \[
    \sqrt[n]{n} > \sqrt[n+1]{n+1} \iff n^{n+1} > {(n+1)}^n \iff n > {\left( \frac{n+1}{n} \right)}^n = {\left( 1 + \frac{1}{n} \right)}^n.
  \]
  Esto vale para todo \(n \geq 3\) porque \({(1 + \frac{1}{n})}^n < 3\) para todo \(n\). En efecto, por el binomio de Newton,
  \[
    {\left( 1 + \frac{1}{n} \right)}^n = \sum_{k=0}^{n} \binom{n}{k} \frac{1}{n^k} \leq \sum_{k=0}^{n} \frac{1}{k!} \leq 1 + \sum_{k=1}^{n} \frac{1}{2^{k-1}} < 1 + 2 = 3,
  \]
  pues \(\binom{n}{k} \frac{1}{n^k} = \frac{n (n-1) \cdots (n-k+1)}{n^k} \cdot \frac{1}{k!} \leq \frac{1}{k!}\) y \(k! \geq 2^{k-1}\).

  Además \(\sqrt[n]{n} \geq 1\) para todo \(n\), de modo que la sucesión converge a un \(l \geq 1\); en particular \(l \neq 0\). Considerando la subsucesión \({(2n)}^{\frac{1}{2n}}\),
  \[
    l^2 = {\left( \lim_{n \to \infty} {(2n)}^{\frac{1}{2n}} \right)}^2 = \lim_{n \to \infty} {(2n)}^{\frac{1}{n}} = \lim_{n \to \infty} 2^{\frac{1}{n}} \cdot \lim_{n \to \infty} n^{\frac{1}{n}} = 1 \cdot l,
  \]
  donde usamos el ejemplo anterior. Luego \(l^2 = l\) y como \(l \neq 0\), \(l = 1\).
\end{proof}

\begin{note}
  La definición de límite puede reformularse así: \(a\) es el límite de \({(x_n)}_{n \in \N}\) si y sólo si para todo \(\e > 0\) el conjunto \(\{ n \in \N : x_n \in (a - \e, a + \e) \}\) tiene complemento finito en \(\N\).
\end{note}

\begin{theorem}
  Un número \(a \in \R\) es el límite de alguna subsucesión de \({(x_n)}_{n \in \N}\) si y sólo si para todo \(\e > 0\) existen infinitos índices \(n\) tales que \(x_n \in (a - \e, a + \e)\).
\end{theorem}

\begin{proof}
  (\(\Rightarrow\)) Sea \(a = \lim_{i \to \infty} x_{n_i}\). Dado \(\e > 0\), existe \(i_0\) tal que \(x_{n_i} \in (a - \e, a + \e)\) para todo \(i > i_0\). Como los \(n_i\) son todos distintos, hay infinitos índices \(n\) con \(x_n \in (a - \e, a + \e)\).

  (\(\Leftarrow\)) Construimos recursivamente índices \(n_1 < n_2 < \cdots\) con \(|x_{n_i} - a| < \frac{1}{i}\). Como \(\{ n : x_n \in (a - 1, a + 1) \}\) es infinito, existe \(n_1\) en él. Dados \(n_1 < \cdots < n_i\), como \(\{ n : x_n \in (a - \frac{1}{i+1}, a + \frac{1}{i+1}) \}\) es infinito, contiene algún \(n_{i+1} > n_i\). Entonces \(|x_{n_i} - a| < \frac{1}{i}\) para todo \(i\), y por lo tanto \(x_{n_i} \to a\).
\end{proof}

\section{Punto de acumulación}

\begin{definition}[Punto de acumulación de una sucesión]
  Un número \(a \in \R\) es un \emph{punto de acumulación} de \({(x_n)}_{n \in \N}\) si es el límite de alguna subsucesión de \({(x_n)}_{n \in \N}\).
\end{definition}

\begin{definition}[Límite inferior y superior]
  Sea \({(x_n)}_{n \in \N}\) una sucesión acotada, digamos \(\alpha \leq x_n \leq \beta\) para todo \(n \in \N\). Sea \(X_n = \{ x_n, x_{n+1}, \dots \}\), de modo que \([\alpha, \beta] \supset X_1 \supset X_2 \supset \cdots\). Llamando \(a_n = \inf X_n\) y \(b_n = \sup X_n\), se tiene
  \[
    \alpha \leq a_1 \leq a_2 \leq \cdots \leq a_n \leq b_n \leq \cdots \leq b_2 \leq b_1 \leq \beta.
  \]
  Por lo tanto existen
  \begin{align*}
    a & = \lim_{n \to \infty} a_n = \sup_{n} \, \inf_{k \geq n} x_k, \\
    b & = \lim_{n \to \infty} b_n = \inf_{n} \, \sup_{k \geq n} x_k.
  \end{align*}
  Se llama \emph{límite inferior} de \({(x_n)}\) a \(a = \liminf_{n \to \infty} x_n\) y \emph{límite superior} a \(b = \limsup_{n \to \infty} x_n\).
\end{definition}

\begin{eg}
  Sea \({(x_n)} = (-1, 2, -\frac{1}{2}, \frac{3}{2}, -\frac{1}{3}, \frac{4}{3}, \dots)\), es decir, \(x_{2n-1} = -\frac{1}{n}\) y \(x_{2n} = 1 + \frac{1}{n}\). Entonces
  \begin{align*}
    \inf X_{2n-2} = \inf X_{2n-1} & = -\frac{1}{n},    \\
    \sup X_{2n-1} = \sup X_{2n}   & = 1 + \frac{1}{n},
  \end{align*}
  de donde \(\liminf_{n \to \infty} x_n = 0\) y \(\limsup_{n \to \infty} x_n = 1\).
\end{eg}
